\documentclass[10pt]{article}
\usepackage[a4paper,margin=23mm]{geometry}
\usepackage[hidelinks]{hyperref}
\usepackage{enumitem}
\usepackage{xcolor}
\setlength{\parindent}{0pt}
\setlength{\parskip}{6pt}
\setlength{\emergencystretch}{2em}
\newcommand{\review}[1]{\textbf{Reviewer comment.} \emph{#1}}
\newcommand{\response}{\textbf{Response.} }

\begin{document}
\begin{center}
{\Large\bfseries Response to Review}\\[5pt]
{\bfseries When does a spectral prior help graph learning?\\
Connectivity-loss estimation under road-network disruptions}\\[5pt]
Van-Truong Le
\end{center}

We thank the reviewer for the careful and constructive assessment. We revised
the manuscript to sharpen its network-science contribution, distinguish
descriptive trends from inferentially supported comparisons, add transport
robustness and size controls, and improve presentation. Locations below refer
to the revised manuscript; source-line references are included to make the
changes easy to audit before journal line numbering is applied.

\section*{Major comments}

\review{M1 --- The contribution relative to residual-learning literature is
not sufficiently explicit.}

\response We agree. Section 2.3 now directly states that bounded residual
learning itself is not claimed as novel and contrasts ResiliRoad with generic
hybrid/physics-guided residual modelling. It identifies the network-science
contribution as the controlled separation of (i) an explicit Fiedler
perturbation prior, (ii) a Fiedler node feature, and (iii) a learned bounded
finite-perturbation correction, followed by geographic transfer and correction
calibration tests. We also added the closely related learnable-physics-engine
literature and explain why the correction-to-prior calibration question is
specific to spectral network disruption. See Section 2.3, p.~4
(source lines 134--159), and the bounded formulation in Section 3.

\review{M2 --- Six country clusters provide limited inferential strength for
leave-one-country-out analysis.}

\response We now separate inference from description explicitly. The opening
paragraph of the Discussion reports that four of the 12 primary expanded-OSM
paired comparisons have clustered 95\% intervals excluding zero and treats the
remaining eight only as descriptive trends. The manuscript therefore does not
claim population-level geographic generality from six countries. The
Limitations section retains the small-cluster warning and states that larger
pre-registered geographic samples are required. See Discussion and
Limitations, p.~16 (source lines 604--643). The statistical unit and nested
bootstrap are defined in Section 4.3 (source lines 333--354).

\review{M3 --- The experiments omit established transportation-robustness
comparators.}

\response We added two non-spectral transport-topology controls: largest
connected-component retention and origin--destination global-efficiency
retention. Section 5.7 reports their association with the algebraic-
connectivity-loss target and explains their intended role. They are useful
controls for whether standard connectivity/reachability summaries explain the
target, but are not presented as capacity- or demand-aware transport models.
This avoids claiming that algebraic connectivity is a universally superior
transport-service metric. See Section 5.7, p.~13 (source lines 489--509), with
the full diagnostic in Appendix Figure 7.

\review{M4 --- Synthetic-to-real transfer confounds graph size and street
morphology.}

\response We added a size-only control spanning 50--1,200 nodes within the same
synthetic graph family. Section 5.7 now begins by stating that this is a
within-family sensitivity analysis, not a morphology-matched transfer test.
The result tests whether size alone reproduces the observed reversal, while the
remaining inability to isolate morphology is retained as a limitation. See
Section 5.7, p.~13 (source lines 489--509), and Limitations, p.~16.

\section*{Requested final checks and presentation revisions}

\review{Check the Journal of Complex Networks length and figure/table rules and
reduce the runtime material if necessary.}

\response We rechecked the current Instructions to Authors. They impose a
300-word abstract maximum and allow up to six keywords, but state no word-count
or figure/table maximum for a Research Article. The revised manuscript has a
sub-300-word abstract, six keywords, approximately 5,000 words by TeXcount, and
seven figures and seven tables including the appendix. Nevertheless, we
condensed Section 5.9 and retained only the runtime evidence needed to support
the screening/scaling claim.

\review{Consolidate repeated statistical hedging.}

\response Done. A single quantitative interpretive rule now opens the
Discussion (p.~16): four of 12 primary expanded-OSM comparisons have clustered
intervals excluding zero; the remaining eight are descriptive. Repetitive
qualifiers were removed from the subsequent Discussion and Conclusion while
the methods and result tables continue to report uncertainty transparently.

\review{Check float numbering and cross-references.}

\response Done. Every figure and table is anchored at its first intended
textual location, and all references use LaTeX \verb|\label|/\verb|\ref| pairs.
The final build has no undefined references, duplicate labels, LaTeX errors, or
overfull boxes. We also moved the architecture-fairness tables to the appendix,
defined ``biased anchor'' at first use (Introduction, source lines 81--84), and
increased axis-label sizes in the multi-panel diagnostics.

\section*{Closing}

No additional experiment is claimed to fully separate road morphology from
other aspects of domain shift. We have instead bounded that conclusion and
made the evidence hierarchy explicit. We believe the revised manuscript now
states its novelty, empirical support, and limitations more precisely while
remaining focused on its central network-science question.

\end{document}
